\documentclass[10pt,a4paper]{amsart}
\usepackage[margin=19mm]{geometry}
\usepackage{amsmath,amssymb,mathtools}
\usepackage[scr=rsfs,cal=euler]{mathalfa}
\usepackage{xfrac}
\usepackage[unicode,hidelinks,bookmarks=false]{hyperref}
\newcommand{\ie}{i.e.\ }
\newcommand{\cf}{cf.\ }
\newcommand{\resp}{respectively\ }
\newcommand{\Subsection}{\subsection}
\providecommand{\nobreakdash}{\nobreak\hbox{-}}
\setlength{\emergencystretch}{2em}
\multlinegap0pt
\usepackage{bm}
\usepackage{bbm}
\usepackage{dsfont}
\usepackage{xspace}
\usepackage{cancel}
\usepackage{mathtools}
\usepackage[shortlabels]{enumitem}
\usepackage{amsthm}
\makeatletter
\@ifundefined{corollary}{\newtheorem{corollary}{Corollary}}{}
\@ifundefined{lemma}{\newtheorem{lemma}{Lemma}}{}
\@ifundefined{proposition}{\newtheorem{proposition}{Proposition}}{}
\@ifundefined{theorem}{\newtheorem{theorem}{Theorem}}{}
\makeatother
\title[Three-level qualitative classification of financial risks un]{Three-level qualitative classification of financial risks under varying conditions through first passage times}
\author{Carlos Bouthelier-Madre, Carlos Escudero}
\date{Source: arXiv:2507.08101v1 (CC BY 4.0)}
\begin{document}
\maketitle

\noindent\emph{Source and licence.} Three-level qualitative classification of financial risks under varying conditions through first passage times. Version arXiv:2507.08101v1 (CC BY 4.0), distributed under the Creative Commons Attribution 4.0 International licence, as recorded in the arXiv licence metadata for this exact version and shown on the abstract page https://arxiv.org/abs/2507.08101v1. Full rights notice: licenses/arxiv-2507.08101-CC-BY-4.0.txt.\par\smallskip

\noindent\textit{Editorial scope.} Counted unit: the Theorem whose source label is \texttt{theor:second\_fpt} in the article (source0.tex lines 720--743, source label \texttt{theor:second\_fpt}), delivered together with the article's own complete original proof of it (source0.tex lines 745--844, one \texttt{proof} environment of 100 lines). No mathematics is written, repaired or re-derived: statement and proof are the article's own text line for line. Required context retained uncounted: lemma:rem2barriers (lines 209--216); mainfpt (lines 224--259); cor:mainfpt (lines 377--399); prop:tlti (lines 426--441); cor:nfpt (lines 637--643). Every definition, display and label of those lines is retained unchanged. Omitted from this package and not redistributed: 1--208, 217--223, 260--376, 400--425, 442--636, 644--719, 744--744, 845--1414. Nothing inside a retained excerpt is omitted or altered: every retained body is delivered exactly as the article's lines are written, so no omission receipt of any kind accompanies this package. Citations: the retained excerpts carry no citation command at all, so no bibliography entry of the article is transcribed and no reference is redistributed. Reference adaptation: none was needed and none was made; every reference inside the delivered unit resolves inside it, so no unresolved reference or citation is produced. Printed slips kept literal: no printed word, symbol, hypothesis or number of the counted statement or of its proof was altered. Layout and class: the article's own class is \texttt{article} with packages including soul, graphicx, color, amsfonts, amssymb, amsthm, amsmath, mathtools, tabularx, comment, cancel, bigints, ulem, scalerel; the wrapper is the lane's pinned \texttt{amsart} editorial wrapper at 10pt on A4, which restates the theorem-like environments the retained text needs and adds the article's own macro definitions that the retained text needs (none), with extra wrapper packages (bm, bbm, dsfont, xspace, cancel, mathtools, enumitem). The delivered numbering is the unit's own. Assets: no retained span contains a figure environment, a graphics-inclusion command, a PDF-page inclusion, a table or a TikZ picture; no image file is copied into this package, the package declares no asset dependency and no image file is redistributed with it. Licence: version arXiv:2507.08101v1 of this article is distributed under the Creative Commons Attribution 4.0 International licence, as recorded in the arXiv licence metadata for this exact version and shown on the abstract page https://arxiv.org/abs/2507.08101v1. A full rights notice is delivered at licenses/arxiv-2507.08101-CC-BY-4.0.txt. Characters: every retained excerpt is pure ASCII, so the delivered engine, XeLaTeX with its default Latin Modern text font, reports no missing character.
\begin{lemma}
\label{lemma:rem2barriers}
Let $\tau_i$ be the FPT of an almost surely continuous stochastic process $C_t$ through a barrier given by the continuous function $B_i(t)$, $i=1,2$, i.e. 
$$
\tau_i = \inf \{t \ge 0 : C_t = B_i(t)\}.
$$
If $B_1(t) \geq B_2(t)\;\forall \, t \geq 0$ and $C_0 > B_1(0)$, then $\tau_1 \leq \tau_2$ almost surely and $\mathbb{E}(\tau_1) \leq \mathbb{E}(\tau_2)$ (be they finite or not). 
\end{lemma}

\begin{theorem}\label{mainfpt}
Consider the stochastic differential equation
$$
dV_t=\mu V_t dt + \sigma V_t dW_t, \qquad \left. V_t \right|_{t=0}= V_0,
$$
with $\mu \in \mathbb{R}$ and $\sigma, V_0>0$. Let the function
\begin{eqnarray}\nonumber
B:[0,\infty) &\longrightarrow& (0,\infty) \\ \nonumber
t &\longmapsto& B(t),
\end{eqnarray}
be such that $B(t) \in \mathcal{C}([0,\infty))$, $K \equiv B(0) < V_0$, and $B(t)>0$ for all $t \ge 0$. Then, if
$$
\tau := \inf \{t \ge 0 : V_t = B(t)\},
$$
it holds that:
\begin{enumerate}
\item[(a)] Whenever either
$$
\liminf_{t \to \infty} B(t) \Big/ \left\{ K \exp \left[(\mu-\sigma^2/2)t - \sigma \sqrt{2t \ln(\ln(t))}\right] \right\} >1,
$$
or
$$
\limsup_{t \to \infty} B(t) \Big/ \left\{ K \exp \left[(\mu-\sigma^2/2)t + \sigma \sqrt{2t \ln(\ln(t))}\right] \right\} >1,
$$
then $\tau < \infty$ almost surely.
\item[(b)] Whenever
$$
B(t) \ge K \exp \left[(\mu-\sigma^2/2)t + \alpha \sqrt{t} \, \right]\;\forall \, t > 0,
$$
for some $\alpha > \sigma$, then $\mathbb{E}(\tau)<\infty$; while if
$$
B(t) \le K \exp \left[(\mu-\sigma^2/2)t + \sigma \sqrt{t} \, \right]\;\forall \, t > 0, 
$$
then $\mathbb{E}(\tau)=\infty$.
\end{enumerate}
\end{theorem}

\begin{corollary}
\label{cor:mainfpt}
The mean value of the stopping time
$$
\tau = \inf \{t \ge 0 : V_t = B(t)\}
$$
through the barrier
$$
B(t) = K \exp \left[(\mu-\sigma^2/2)t + \alpha \sqrt{t} \, \right],
$$
for any $\alpha > \sigma$, is given by the explicit formula:
$$
\mathbb{E}(\tau) = \frac{(\ln(V_0/K))^2}{\alpha^2 - \sigma^2}.
$$
Moreover, the mean FPT for a barrier that fulfills
$$
B(t) \ge K \exp \left[(\mu-\sigma^2/2)t + \alpha \sqrt{t} \, \right],
$$
for all $t > 0$ and some $\alpha > \sigma$, obeys the explicit bound
$$
\mathbb{E}(\tau) \le \frac{(\ln(V_0/K))^2}{\alpha^2 - \sigma^2}.
$$
\end{corollary}

\begin{proposition}\label{prop:tlti}
Let $V_t$ be the unique solution to the SDE
$$
dV_t=\mu V_t dt + \sigma V_t dW_t, \qquad \left. V_t \right|_{t=0}= V_0,
$$
with $\mu \in \mathbb{R}$ and $\sigma, V_0>0$, and the barrier function
\begin{eqnarray}\nonumber
B:[0,\infty) &\longrightarrow& (0,\infty) \\ \nonumber
t &\longmapsto& B(t),
\end{eqnarray}
be such that $B(t) \in \mathcal{C}([0,\infty))$ and $K \equiv B(0)<V_0$, but arbitrary otherwise. Then, for the FPT
$$
\tau := \inf \{t \ge 0 : V_t = B(t)\},
$$
it holds that $\mathbb{P}(\tau < \infty) >0$.
\end{proposition}

\begin{corollary}\label{cor:nfpt}
Let $V_t$ and $B(t)$ be as in the statement of Proposition~\ref{prop:tlti}. If we define the FPT as
$$
\tau := \inf \{0 \le t \le T : V_t = B(t)\}
$$
with the convention $\inf \emptyset = \infty$, then $\mathbb{P}(\tau < \infty) < 1$ for any $T>0$.
\end{corollary}

\begin{theorem} 
\label{theor:second_fpt}
Denote $B_c(t,\alpha):=K \exp \left[(\mu-\sigma^2/2)t + \alpha \sqrt{t} \, \right]$ and let $B(t)$ be a barrier function that fulfills the same properties as those stated in Theorem~\ref{mainfpt}. Define $\tau := \inf \{t \ge 0 : V_t = B(t)\}$; then:
\begin{enumerate}[label=(\alph*)]
\item If, for some $\alpha>\sigma$,
$$
\liminf\limits_{t\rightarrow\infty}\dfrac{B(t)}{B_c(t,\alpha)} > 1,
$$
 then $\mathbb{E}(\tau)<\infty$. In particular, with the Gaussian probability distribution function denoted by
$$
\phi_{a,b}(x):=\frac{1}{\sqrt{2\pi} \, b}\exp \left\lbrace -\frac{(x-a)^2}{2 \, b^2} \right\rbrace,
$$
the following upper bound holds
\begin{equation}\nonumber
\mathbb{E}(\tau) \le T + \int_{\alpha\sqrt{T}}^\infty \frac{x^2}{\alpha^2 - \sigma^2} \, \phi_{q,\sigma\sqrt{T}}(x) \, dx,
\end{equation}
where $T:=\inf \left\lbrace t\geq 0 \, \big\vert \, B(s) \ge B_c(s,\alpha) \; \forall \, s \ge t \right\rbrace$ and $q:=\ln(V_0/K)$.
\item If, conversely, 
$$
\limsup\limits_{t\rightarrow\infty}\dfrac{B(t)}{B_c(t,\sigma)} < 1,
$$
then $E(\tau)=\infty$.
\end{enumerate}
\end{theorem}

\begin{proof}
For part \textit{(a)} of the statement, note $\tau<\infty$ almost surely by part \textit{(a)} of Theorem~\ref{mainfpt}. Also, by the continuity of $B(\cdot),B_c(\cdot,\alpha)$, their relative limit behavior, and the intermediate value theorem, $\exists \, T \ge 0$ such that:
$$
T:=\inf \left\lbrace t\geq 0 \, \big\vert \, B(s) \geq B_c(s,\alpha) \; \forall \, s \ge t \right\rbrace<\infty.
$$
This deterministic time denotes the last crossing point of the graphs of $B(\cdot)$ and $B_c(\cdot,\alpha)$ provided it exists, being zero otherwise; by crossing we mean that the relative order of these graphs changes at this point. If $T=0$, this brings us back to the cases analyzed in Theorem~\ref{mainfpt} and Corollary~\ref{cor:mainfpt}, so from now on we assume $T>0$. By the proof of part~\textit{(a)} of Theorem~\ref{mainfpt}, we might rewrite the first passage time as:
\begin{eqnarray}\nonumber
\tau &=& \inf \{t \ge 0 : \ln(V_t) = \ln(B(t))\}\\ \nonumber
&=& \inf \{t \ge 0 : W_t = \ln(B(t))/\sigma - \left[(\mu - \sigma^2/2)t + \ln(V_0) \right]/\sigma \}.
\end{eqnarray}
Now, by the law of the total expectation,
\begin{eqnarray}
\label{eq:total_expectation}
\nonumber
\mathbb{E}(\tau)&=&\mathbb{E}(\tau 1_{\tau\leq T})+\mathbb{E}(\tau 1_{\tau> T})\\ \nonumber
&\leq& T \, \mathbb{P}(\tau\leq T)+\mathbb{E}(\tau 1_{\tau> T}) \\
&=& T+\mathbb{E}[(\tau-T) 1_{\tau> T}],
\end{eqnarray}
and for the last term, again by this law, we can compute:
$$
\mathbb{E}[(\tau-T) 1_{\tau> T}]=\int\limits_{B(T)}^\infty\mathbb{E}(\tau-T\vert\tau>T,V_T=x) \, \mathbb{P}(V_T\in dx, \tau>T),
$$
where the probability density can be decomposed as the product
$$\mathbb{P}(V_T\in dx, \tau>T)=\mathbb{P}(V_T\in dx) \, \mathbb{P}( \tau>T\vert V_T =x).
$$
For brevity, we will denote $L_T(x) dx:=\mathbb{P}(V_T\in dx)$, referencing the lognormal distribution of $V_T$.

We will now focus on $\mathbb{E}(\tau-T\vert\tau>T,V_T=x)$. 
%Let us define $\tilde{B}(t)=\ln\left(\dfrac{B(t)}{K \exp[(\mu-\sigma/2)t]}\right)$
Firstly, consider the random time defined as $\tau^\prime:=(\tau-T)1_{\tau>T}$. That is, $\tau^\prime=0$ if $\tau \leq T$ and else it is the stopping time
\begin{eqnarray}\nonumber
\tau' &=& \inf \{ t \ge 0 : V_{t+T}=B(t+T)\}\\ \nonumber
&=&  \inf \{ t \ge 0 : W_{t+T}-W_{T} = \ln[B(t+T)/V_T]/\sigma - \left[(\mu - \sigma^2/2)t \right]/\sigma \}\\ \nonumber
&=&  \inf \{ t \ge 0 : \sigma(W_{t+T}-W_{T})+q+\sigma W_T = \tilde{B}(t+T)\},
\end{eqnarray}
where we have used that $\tilde{B}(t):=\ln\left(B(t)/[K \exp\{(\mu-\sigma^2/2)t\}]\right)$ and $W_T=[\ln(V_T/V_0)-(\mu-\sigma^2/2)T]/\sigma$. Thus, for each $x>B(T)$ (note that $V_T>B(T)\Leftrightarrow \sigma W_T +q> \alpha\sqrt{T}$, by the definition of $T$) we have:
$$
\mathbb{E}(\tau-T\vert V_T=x,\tau>T) =\mathbb{E}(\tau^\prime\vert V_T=x,\tau>T)=\mathbb{E}(\tau^\prime\vert V_T=x),
$$ 
where the last equality is obtained by the Markovian nature of $W_t$. Besides, $B(t+T)\geq K \exp\{(\mu-\sigma^2/2)(t+T)+\alpha\sqrt{t+T}\}$ by the definition of $T$, implying
\begin{eqnarray}\nonumber
B(t+T) &\geq& K \exp\{(\mu-\sigma^2/2)T\} \\ \nonumber
&& \times
\exp\{(\mu-\sigma^2/2)t+\alpha \sqrt{t+T}\} \\ \nonumber
 &>& K \exp\{(\mu-\sigma^2/2)T\} \\ \nonumber
&& \times
\exp\{(\mu-\sigma^2/2)t+\alpha \sqrt{t}\} \\ \nonumber
&=& B_c(t,\alpha) \exp\{(\mu-\sigma^2/2)T\}.
\end{eqnarray}
By the translation invariance of Brownian motion, the stochastic process
$\{W_{t+T}-W_{T},t \ge 0\}$ is a standard Brownian motion and, by the independence of Brownian increments, it is independent of $W_T$ (and of $V_T$ since $V_T$ is $\sigma(W_T)-$measurable). Therefore, $\tau'$ for a fixed value of $V_T$ (or equivalently $W_T$), falls under the assumptions of Theorem~\ref{mainfpt} and Corollary~\ref{cor:mainfpt}, 
so it is summable and moreover
$$
\mathbb{E}(\tau' \, \vert \, W_T=z) \le \frac{(q +\sigma z)^2}{\alpha^2 - \sigma^2}.
$$
Now, substituting this estimate in equation~\eqref{eq:total_expectation}, we have:
\begin{eqnarray}
\label{eq:for_lemma}
\nonumber
\mathbb{E}(\tau) &\le& T + \int\limits_{B(T)}^\infty\frac{|q+\sigma z(x)|^2}{\alpha^2 - \sigma^2} \, \mathbb{P}(\tau>T\vert V_T=x) \, L_T(x) \, dx \\
&=&
T + \int\limits_{\alpha\sqrt{T}}^\infty\frac{y^2}{\alpha^2 - \sigma^2} \, \mathbb{P}(\tau>T\vert \sigma W_T+q=y) \, \phi_{q,\sigma \sqrt{T}}(y) \, dy,
\end{eqnarray}
where we have used the relations $y=\sigma z(x)+q=\ln(x/K)-(\mu-\sigma^2/2)T$. Note that the change in the integration limits comes from the fact that $V_T>B(T)\Leftrightarrow \sigma W_T+q>\alpha\sqrt{T}$ due to the definition of $T$, which implies $B(T)=B_c(T,\alpha)$. Also, note that, since $V_T$ is lognormally distributed, then $q+\sigma W_T$ is normally distributed with mean $q$ and variance $\sigma^2T$, hence the appearance of $\phi_{q,\sigma\sqrt{T}}(y)$ in the second integrand. To obtain the expression in the statement, simply consider
$$
\mathbb{P}(\tau>T\vert \sigma W_T+q=y)\leq 1
$$
uniformly in $y$.

The proof of part \textit{(b)} is as follows. Firstly, we consider the deterministic crossing time $T_c$ of the barrier $B(t)$ through the critical barrier $B_c(t,\sigma)$. Indeed, by the continuity of $B(\cdot),B_c(\cdot,\sigma)$, their relative limit behavior, and the intermediate value theorem, $\exists \, T_c \ge 0$ such that:
$$
T_c :=\inf \left\lbrace t\geq 0 \, \big\vert \, B(s) \le B_c(s,\sigma) \; \forall \, s \ge t \right\rbrace<\infty.
$$
As before, the case $T_c=0$ is analyzed in Theorem~\ref{mainfpt}, so $T_c>0$ will be assumed. We introduce two random times, $\tau^\prime_c:=(\tau-T_c)1_{\tau>T_c}$ and
$$
\tau_{1}:=\inf \left\lbrace t\geq 0 \, \big\vert \sigma(W_{t+T_c}-W_{T_c})+q+\sigma W_{T_c}=\sigma\sqrt{t+T_c}\right\rbrace.
$$
By the definition of $T_c$, $\tilde{B}(t+T_c)\leq \sigma\sqrt{t+T_c}$. Thus, $\tau^\prime_c \geq \tau_1 \, 1_{\tau>T_c}$ almost surely by Lemma~\ref{lemma:rem2barriers}. We add now another FPT,
$$
\tau_{2}:=\inf \left\lbrace t\geq 0 \, \big\vert \sigma(W_{t+T_c}-W_{T_c})+q+\sigma W_{T_c}=\sigma\left(\sqrt{t}+\sqrt{T_c}\right)\right\rbrace,
$$
and we realize that, by the subadditivity of the square root, i.e. $\sqrt{t+T_c}\leq \sqrt{t}+\sqrt{T_c}\;\forall \, t,T_c\geq 0$, it holds that $\tau_1 \, 1_{\tau>T_c} \geq \tau_2 \, 1_{\tau>T_c}$ almost surely again by Lemma~\ref{lemma:rem2barriers}. As in the first part of the proof, for each fixed $W_{T_c}$ such that $V_{T_c} > B(T_c)$, equivalently $q+\sigma W_{T_c}-\sigma\sqrt{T_c}>0$, $\tau_2$ falls under the assumptions of Theorem~\ref{mainfpt}, but in this case it fulfills the conditions for an infinite mean. Thus, since $\tau^\prime_c \geq \tau_1 \, 1_{\tau>T_c} \geq \tau_2 \, 1_{\tau>T_c}$ almost surely, we find that
$$
\mathbb{E}(\tau^\prime_c\vert V_{T_c}=y,\tau>T_c) \geq \mathbb{E}(\tau_2\vert V_{T_c}=y,\tau>T_c) = \mathbb{E}(\tau_2\vert V_{T_c}=y)= \infty
$$
for every $y>B(T_c)$, where we have employed the Markovianity of Brownian motion.
Using that $1_{\tau>T_c}\leq 1$ and $T_c > 0$, we find
\begin{equation}\nonumber
\mathbb{E}(\tau) \ge \mathbb{E}(\tau 1_{\tau>T_c}) > \mathbb{E}((\tau-T_c) 1_{\tau>T_c}).
\end{equation}
Now, as before, by the law of the total expectation:
\begin{equation}\nonumber
\mathbb{E}(\tau^\prime_c) = \int\limits_{B(T_c)}^\infty \mathbb{E}(\tau^\prime_c \vert V_{T_c}=y,\tau>T_c) \, \mathbb{P}(\tau>T_c\vert V_{T_c}=y) \, \mathbb{P}(V_{T_c} \in dy).
\end{equation}
Since $\mathbb{P}(\tau > T_c) >0$  by Corollary~\ref{cor:nfpt} and
$$
\mathbb{P}(\tau > T_c)=\int\limits_{B(T_c)}^\infty \mathbb{P}(\tau>T_c\vert V_{T_c}=y) \, \mathbb{P}(V_{T_c}\in dy),
$$
we know that $\mathbb{P}(\tau>T_c\vert V_{T_c}=y)$ shares non-empty support with $\mathbb{P}(V_{T_c}\in dy)$ for $y>B(T_c)$. This implies $\mathbb{E}(\tau^\prime_c)=\infty$ and consequently $\mathbb{E}(\tau)=\infty$.
\end{proof}

\end{document}
